compact subintervals of $(0,\infty)$ gives $w(y)=y^5 f^{(4)}(y)/12$.
An elementary differentiation identity is
\[
 \frac{d^4}{dy^4}\frac{y^5}{y^2+a^2}
       =24a^4\operatorname{Im}(a-iy)^{-5}.
\]
Integrate by parts four times first on $[\varepsilon,R]$. At zero the
boundary term involving the $k$th derivative of $y^5/(y^2+a^2)$ and
the $(3-k)$th derivative of $f$, for $0\le k\le3$, is
$O(y^{5-k}y^{-4+k})=O(y)$. At infinity it tends to zero exponentially.
The differentiated integrands are absolutely integrable, so passage to
the endpoints is justified. Hermite's integral formula at $s=5$ now gives
\[
 \int_0^\infty\frac{w(y)}{y^2+a^2}\,dy
 =2a^4\int_0^\infty\frac{\operatorname{Im}(a-iy)^{-5}}
                           {e^{2\pi y}-1}\,dy
 =a^4\zeta(5,a)-\frac1{2a}-\frac14.
\]
At $a=j\in\Z_{>0}$, use $\zeta(5,j)=\xi-\Hfive{j-1}$ and
$\Hfive j-\Hfive{j-1}=j^{-5}$ to recover \eqref{eq:mupole}. This verifies
\eqref{eq:momentrep} on the partial-fraction basis and hence on its domain.

Finally, a nonzero real polynomial $q$ of degree less than $h$ satisfies
\[
 \boldsymbol q^{\mathsf T}G_K(\xi)\boldsymbol q
 =\int_0^\infty\frac{D_N(y^2)^6q(y^2)^2}{D_K(y^2)}w(y)\,dy>0.
\]
The integrand is continuous and positive on some interval: $q$ cannot
vanish on all positive real numbers. This proves positive definiteness
and the determinant assertion.
\end{proof}

\section{The local rational functional: explicit domains and distribution}
\label{sec:local}

\subsection{Polynomial and integer-pole definitions}
Let $L(x^k)=B_k$ on $\Q[x]$, and let $\tau(P)=L(P''')/24$. For $r\in\Z$,
define the nonnegative integer
\[
 d(r)=\begin{cases}r,&r\ge0,\\-r-1,&r<0,\end{cases}\qquad
 \tau_X\left(\frac1{x-r}\right)=\Hfive{d(r)}-X.
\]
Extend by polynomial division and finite simple partial fractions to the
vector space of rational functions with polynomial part and simple integer
poles. At a fixed prime, extend all these rational functionals coefficientwise
from $\Q$ to $\Qp$ before applying them to $p$-adic inputs. The following identities include the required domain restriction:
\begin{align}
 \mu_X(R)&=\tau_X\bigl(x^5R(-x^2)\bigr),\label{eq:pullback}\\
 \tau_X(g(-1-x))&=-\tau_X(g(x)),\label{eq:reflection}\\
 \tau_X(g(x+1)-g(x))&=g^{(4)}(0)/24
                  &&\text{if $g$ is regular at $0$}.\label{eq:difference}
\end{align}
For \eqref{eq:pullback}, monomials follow by differentiating three times.
For a pole use
\[
 \frac{x^5}{j^2-x^2}=-x^3-j^2x-\frac{j^4}{2}
          \left(\frac1{x-j}+\frac1{x+j}\right),
\]
then $d(j)=j$ and $d(-j)=j-1$. The polynomial cases of
\eqref{eq:reflection}--\eqref{eq:difference} follow from Bernoulli
reflection and difference. For the pole cases use $d(-1-r)=d(r)$ and,
for $r\ne0$,
\[
 \Hfive{d(r-1)}-\Hfive{d(r)}=-r^{-5}.
\]
The last identity holds for negative as well as positive $r$.

\subsection{Restricted power series and meromorphic direct sums}
Fix a prime $p\ge7$. Let
\[
 \mathcal A=\Qp\langle z\rangle
 =\{\textstyle\sum_{d\ge0}f_dz^d:\ f_d\in\Qp,\ f_d\to0\},
 \qquad \|f\|=\sup_d|f_d|_p.
\]
Completeness follows by coefficientwise completeness and uniform control
of the tails. Polynomial truncations are dense. Put
\[
 \kappa_d=\frac{d(d-1)(d-2)B_{d-3}}{24}\ (d\ge3),
 \qquad\kappa_0=\kappa_1=\kappa_2=0.
\]
Von Staudt--Clausen implies $v_p(\kappa_d)\ge-1$ for every $d$, and
$\kappa_d\in\Zp$ for $d\le p+1$. Hence
\begin{equation}\label{eq:tauextseries}
 \tau\left(\sum_{d\ge0}f_dz^d\right)=\sum_{d\ge0}f_d\kappa_d,
 \qquad \|\tau\|\le p.
\end{equation}
The series converges because its terms tend to zero in the complete
nonarchimedean field.

For a finite set $S$ of distinct integers, let
\[
 \mathcal B_S=\mathcal A\oplus\bigoplus_{r\in S}\Qp\frac1{z-r}.
\]
Give this space the maximum of the analytic coefficient norm and residue
norms. It embeds into the localization of $\mathcal A$ obtained by
inverting $\prod_{r\in S}(z-r)$. Indeed, if
$f+\sum_rc_r/(z-r)=0$, multiplication by that product and evaluation at
$r$ gives $c_r\prod_{s\ne r}(r-s)=0$, so each residue vanishes and then
$f=0$. This proves uniqueness of the decomposition. No claim that
$\mathcal B_S$ itself is a ring is needed.

For $Y\in\Qp[X]$ define
\begin{equation}\label{eq:taudirectsum}
 \taus_Y\left(f+\sum_{r\in S}\frac{c_r}{z-r}\right)
   =\tau(f)+\sum_{r\in S}c_r(\Hfive{d(r)}-Y).
\end{equation}
These maps are compatible when $S$ is enlarged. Their images lie in the
finite-dimensional space $\Qp+\Qp Y\subset\Qp[X]$, which is complete for
the Gauss norm. In particular, convergence in this proof never relies on
completeness of the entire unbounded-degree polynomial ring.

\begin{lemma}[Bounded division in the restricted power-series space]
\label{lem:boundeddivision}
Let $T(z)=\prod_{\nu=1}^t(z-r_\nu)$, where $r_\nu\in\Z$ and every
$r_\nu-r_\eta$, $\nu\ne\eta$, is a $p$-adic unit. For $f\in\mathcal A$,
\[
 \frac fT=P+\sum_{\nu=1}^t\frac{c_\nu}{z-r_\nu},\qquad
 c_\nu=\frac{f(r_\nu)}{T'(r_\nu)},
\]
with $P\in\mathcal A$ and
$\max(\|P\|,|c_1|_p,\ldots,|c_t|_p)\le\|f\|$.
For a polynomial $f$ this is ordinary division and partial fractions;
its polynomial quotient has degree at most $\deg f$.
\gid{Zeta5.Local.bounded_division}
\end{lemma}
\begin{proof}
Evaluation at $r\in\Zp$ has norm at most one. More explicitly, if
$f=\sum f_kz^k$, then
\[
 \frac{f(z)-f(r)}{z-r}=\sum_{n\ge0}
         \left(\sum_{k\ge n+1}f_kr^{k-1-n}\right)z^n.
\]
The inner series converges, its absolute value is at most
$\sup_{k\ge n+1}|f_k|_p$, and these bounds tend to zero. This proves
bounded division by one linear factor after removing the remainder.

For polynomial $f$, division by the monic integral polynomial $T$ and
repeated subtraction of its leading term give a quotient and remainder
with coefficients bounded by $\|f\|$. Since $T'(r_\nu)$ is a unit, the
partial-fraction residues have the same bound. Apply this polynomial map
to the truncations of an arbitrary $f$ and use completeness of
$\mathcal B_S$. The displayed evaluation formula for $c_\nu$ passes to
the limit. The localization identity passes to the limit after
multiplication by $T$; uniqueness was proved above.
\end{proof}

\begin{proposition}[Integral extension with a degree-restricted leading term;
source Lemma 3.1]
\label{prop:integralextension}
Under the hypotheses of Lemma~\ref{lem:boundeddivision}, suppose also
$d(r_\nu)<p$, $Y\in\Zp[X]$, and $U_j\in\Zp[z]$ for every $j\ge0$.
If $\deg U_0\le p+1$, then the series below converges in $\mathcal B_S$ and
\begin{equation}\label{eq:integralextension}
 \taus_Y\left(\sum_{j\ge0}\frac{p^jU_j(z)}{T(z)}\right)\in\Zp[X].
\end{equation}
\gid{Zeta5.Local.integral_extension}
\end{proposition}
\begin{proof}
Lemma~\ref{lem:boundeddivision} bounds the norm of each $U_j/T$ by one;
the factors $p^j$ force convergence. Each principal-part value
$\Hfive{d(r_\nu)}-Y$ is integral because $d(r_\nu)<p$. The analytic
functional loses at most one factor of $p$, by
\eqref{eq:tauextseries}. For $j=0$, the quotient is a polynomial of degree
at most $p+1$, so its value is integral without that loss. For $j\ge1$,
$p^j$ absorbs it. All partial sums have integral output. Their limit in
$\Qp+\Qp Y$ is integral coefficientwise because $\Zp$ is closed.
The case $T=1$ is included.
\end{proof}

\begin{lemma}[Far poles and order of expansion]
\label{lem:farpoles}
If $s\in\Qp$ has $v_p(s)<0$, then $(z-s)^{-1}\in\mathcal A$ and
\begin{equation}\label{eq:farseries}
 \tauan\left(\frac1{z-s}\right)
  =-\frac14\sum_{k\ge0}\binom{k+3}{3}B_k s^{-k-4}.
\end{equation}
If a rational function has polynomial part and only simple poles, each
pole being either an integer or outside $\Zp$, its value under
\eqref{eq:taudirectsum} is independent of whether one first takes global
partial fractions or expands its far denominator factors in
$\mathcal A$ and then separates its near principal parts.
\gid{Zeta5.Local.far_pole_compatibility}
\end{lemma}
\begin{proof}
The geometric series
$(z-s)^{-1}=-\sum_{d\ge0}s^{-d-1}z^d$ converges in the coefficient norm.
Apply the continuous functional and substitute $d=k+3$ to obtain
\eqref{eq:farseries}. Products of the far expansions converge in
$\mathcal A$ and represent the corresponding rational functions there.
Both proposed procedures therefore give the same element of the
localization of $\mathcal A$ at the near-pole factors. Both have the same
simple principal parts and analytic remainder by uniqueness. Applying the
single map \eqref{eq:taudirectsum} proves the assertion. No use of the
integer-pole harmonic formula at a nonintegral far pole is permitted.
\end{proof}

\begin{lemma}[Reflection and difference on the completed domain]
\label{lem:completedidentities}
On the union of the spaces $\mathcal B_S$, for finite integer pole sets
$S$, reflection satisfies
\[
 \taus_Y(g(-1-z))=-\taus_Y(g(z)).
\]
If $g$ is regular at zero, translation satisfies
\[
 \taus_Y(g(z+1)-g(z))=g^{(4)}(0)/24.
\]
Both assertions include analytic rational functions with far poles via
Lemma~\ref{lem:farpoles}.
\gid{Zeta5.Local.completed_translation_reflection}
\end{lemma}
\begin{proof}
The substitutions $z\mapsto z+1$ and $z\mapsto-1-z$ are bounded maps of
$\mathcal A$ into itself: the binomial coefficients are integral and the
polynomial substitution maps have norm at most one, so they extend from
dense polynomial truncations. Differentiation and evaluation at zero are
also continuous. Thus the polynomial identities
\eqref{eq:reflection}--\eqref{eq:difference} extend to the analytic part.
For an integer principal part reflection sends its pole to $-1-r$, and
translation sends it to $r-1$. The harmonic identities already checked
prove the result on these terms. In the difference identity the term
$r=0$ is excluded exactly by regularity at zero. Linearity combines the
parts. A far pole remains far under these unit translations and
reflection, so its analytic interpretation is consistent throughout.
\end{proof}

\subsection{Distribution, with the pole recurrence made explicit}
Define
\begin{equation}\label{eq:CpY}
 C_p=\sum_{a=1}^{p-1}\tauan\left(\frac1{z+a/p}\right),
 \qquad Y=p^5X+C_p.
\end{equation}

\begin{lemma}[Integrality of the distribution parameter]
\label{lem:Cp}
One has $C_p\in p^5\Zp$, and hence $Y\in\Zp[X]$.
\gid{Zeta5.Local.distribution_parameter_integral}
\end{lemma}
\begin{proof}
The $k=0$ terms of \eqref{eq:farseries} sum to
$-p^4\sum_{a=1}^{p-1}a^{-4}/4$. Modulo $p$ this sum of fourth inverse
powers is zero: $p-1\nmid4$ for $p\ge7$, and the multiplicative group of
$\mathbb F_p$ gives the standard power-sum identity. The $k=1$ terms
have valuation at least $5$. Each $k\ge2$ term has valuation at least
$k+3\ge5$, using $v_p(B_k)\ge-1$. The series converges, so the claimed
closed ideal contains its sum.
\end{proof}

\begin{proposition}[Full distribution identity; source Lemma 3.2]
\label{prop:distribution}
Let $g\in\Qp(z)$ have polynomial part and finitely many simple integer
poles. With $Y$ as in \eqref{eq:CpY},
\begin{equation}\label{eq:distribution}
 \tau_X(g)=p^{-4}\sum_{a=0}^{p-1}\taus_Y\bigl(g(a+pz)\bigr)
 \quad\text{in }\Qp[X].
\end{equation}
In each summand a pole $(r-a)/p$ is an integer when $r\equiv a\pmod p$,
and has valuation $-1$ otherwise; it is interpreted accordingly. No
restriction $d(r)<p$ is imposed on the original integer poles in this
proposition. That restriction belongs only to the integrality estimate.
\gid{Zeta5.Local.distribution}
\end{proposition}
\begin{proof}
All summands are in the completed domain by
Lemma~\ref{lem:farpoles}. Denote the right side by $T(g)$.
For a polynomial, Bernoulli multiplication gives
$\sum_{a=0}^{p-1}L(P(a+pz))=pL(P)$.
Three differentiations contribute a further factor $p^3$, hence
$T(P)=\tau(P)$. These identities are rational polynomial identities and
therefore hold for coefficients in $\Qp$ as well.

For $g_0(z)=1/z$, definition \eqref{eq:CpY} gives
\[
 T(g_0)=p^{-5}(-Y+C_p)=-X.
\]
Next, replacing $a$ by $p-1-a$ and $z$ by $-1-z$ in the sum, and applying
Lemma~\ref{lem:completedidentities}, gives
$T(g(-1-z))=-T(g(z))$. If $g$ is regular at zero, telescoping the finite
sum gives
\begin{align*}
 T(g(z+1)-g(z))
 &=p^{-4}\taus_Y\bigl(g(p(z+1))-g(pz)\bigr)\\
 &=p^{-4}\frac{(g(pz))^{(4)}|_{z=0}}{24}
 =\frac{g^{(4)}(0)}{24}.
\end{align*}
Here $g(pz)$ belongs to the completed domain and remains regular at zero;
thus the completed difference identity applies with precisely its stated
hypothesis.

For $r>0$ let $g_r(z)=1/(z-r)$. Since $g_r$ is regular at zero,
\[
 T(g_{r-1})-T(g_r)=-r^{-5}.
\]
Induction from $T(g_0)=-X$ yields $T(g_r)=\Hfive r-X$ for every $r\ge0$.
Reflection gives $T(g_{-1-r})=T(g_r)$ and hence all negative integer
poles. Agreement on polynomials and simple principal parts proves
\eqref{eq:distribution}. This argument is independent of any truncated
$p$-adic numerical test.
\end{proof}

